\section{Path composition}\label{sec:11-path-algebras:05-path-composition:1}

Given a path \texttt{p : Path Q u v} and a path \texttt{q : Path Q v w}, they can be concatenated into a path \texttt{Path Q u w}, defined by recursion on \texttt{q}:

\begin{lstlisting}
def Path.append {V A : Type} {Q : Quiver V A} {u v w : V}
    (p : Path Q u v) (q : Path Q v w) : Path Q u w :=
  match q with
  | Path.nil _ => p
  | Path.cons a h h' q' => Path.cons a h h' (Path.append p q')
\end{lstlisting}

Reading the recursion:

\begin{itemize}
\tightlist
\item
  If \texttt{q} is trivial (\texttt{Path.nil \_\allowbreak{}}, a path of length zero from \texttt{v} to \texttt{v}), then appending \texttt{p} before it just gives back \texttt{p} itself (there is nothing to append).
\item
  If \texttt{q} ends with an arrow \texttt{a} (i.e.~\texttt{q = Path.cons a h h' q'} for some shorter path \texttt{q'}), then appending \texttt{p} before all of \texttt{q} is the same as appending \texttt{p} before the shorter \texttt{q'}, and then re-attaching the same final arrow \texttt{a} on top.
\end{itemize}

This recursion terminates because each recursive call is on the strictly shorter path \texttt{q'}. This is the same structural recursion principle behind the \texttt{induction} tactic from Chapter 4.

\subsection{\texorpdfstring{A worked instance: rebuilding \texttt{pathBetaAlpha} via \texttt{append}}{A worked instance: rebuilding pathBetaAlpha via append}}\label{sec:11-path-algebras:05-path-composition:2}

The previous section built \texttt{pathBetaAlpha : Path exampleQuiver 0 2} directly, one \texttt{Path.cons} at a time. Here it is again, built instead by \emph{composing} the shorter path \texttt{pathAlpha} with a fresh one-arrow path using \texttt{Path.append}. Both constructions produce the exact same path.

\begin{lstlisting}
def pathBetaOnly : Path exampleQuiver 1 2 :=
  Path.cons ExampleArrow.beta rfl rfl (Path.nil 1)
\end{lstlisting}

\texttt{pathBetaOnly} is that fresh one-arrow path: a path from \texttt{1} to \texttt{2} consisting of just the single arrow \texttt{beta}, built on its own rather than by extending \texttt{pathAlpha}.

\begin{lstlisting}
def pathBetaAlphaViaAppend : Path exampleQuiver 0 2 :=
  Path.append pathAlpha pathBetaOnly
\end{lstlisting}

\texttt{pathBetaAlphaViaAppend} composes the shorter path \texttt{pathAlpha} with \texttt{pathBetaOnly} via \texttt{Path.append}, yielding a path from \texttt{0} to \texttt{2}. This has the same endpoints as \texttt{pathBetaAlpha}, but is assembled by composition rather than by chaining \texttt{Path.cons} calls directly.

\begin{lstlisting}
example : pathBetaAlphaViaAppend = pathBetaAlpha := rfl
\end{lstlisting}

That final \texttt{rfl} is not a weak or approximate check. It says the two constructions are \emph{definitionally} the same term, reducing to an identical normal form (Chapter 5), not merely ``provably equal after some argument.'' This is the concrete payoff of \texttt{Path.append}'s recursive definition: composing paths built independently, via arbitrarily different routes, agrees exactly with building the composite path directly by hand, because both ultimately unfold to the same sequence of \texttt{Path.cons} applications.

\begin{mathreading}

\texttt{Path.append} is \textbf{composition in the free category} \(\mathrm{Free}(Q)\), written throughout this section in \emph{path order}: ``\(p\) then \(q\),'' matching the argument order of \texttt{Path.append p q}. This is not the function-composition order \(q \circ p\) standard in category theory texts. (The two conventions denote the same composite; only the order in which the symbols are written differs, and mixing them mid-explanation is a common source of confusion, so this book fixes path order throughout.) In path order, \texttt{Path.append} is the map

\end{mathreading}

\[
\mathrm{Hom}(u,v) \times \mathrm{Hom}(v,w) \longrightarrow \mathrm{Hom}(u,w),
\qquad (p, q) \longmapsto p\,;\,q
\]

(the semicolon ``\(;\)'' is a common notation for path-order composition, distinguishing it from function-order \(\circ\)). The identity laws, stated in path order: \texttt{Path.append p (Path.nil v) = p} (appending the trivial path \emph{after} \texttt{p} changes nothing; this case is implemented directly by the \texttt{nil} branch of the \texttt{match} above, since recursion is on \texttt{q}), and, separately, \texttt{Path.append (Path.nil u) p = p} (appending the trivial path \emph{before} \texttt{p} also changes nothing; proved as Exercise 2 by induction on \texttt{p}, since the \texttt{nil} branch alone does not give this one for free). The \texttt{cons} case of the recursion is exactly associativity of concatenation. Together with \texttt{nil} as identities, this makes \(\mathrm{Free}(Q)\) a genuine category: the smallest/most general category containing \(Q\)'s arrows, in the sense of a \hyperref[sec:01-basics:04-terminology:1]{universal property}.

\textbf{Mathlib equivalent.} \texttt{Path.append} is already in Mathlib, as \texttt{Quiver.Path.comp}. It is the same recursion (on the \emph{second} path), with the same ``\texttt{nil} does nothing, \texttt{cons} re-attaches its trailing arrow'' shape:

\begin{lstlisting}
def pathBetaOnly' : Quiver.Path (1 : Fin 3) 2 := Quiver.Path.cons Quiver.Path.nil MyArrow.beta
def pathBetaAlphaViaComp' : Quiver.Path (0 : Fin 3) 2 := Quiver.Path.comp pathAlpha' pathBetaOnly'

example : pathBetaAlphaViaComp' = pathBetaAlpha' := rfl
\end{lstlisting}

This is the same closing \texttt{rfl} as the book's \texttt{pathBetaAlphaViaAppend = pathBetaAlpha} check: two paths built via different routes (direct \texttt{cons}-chaining versus composing two shorter paths) reduce to the identical term, because \texttt{Quiver.Path.comp} unfolds to exactly the same sequence of \texttt{cons} applications \texttt{Path.append} does.

\subsection{The path algebra}\label{sec:11-path-algebras:05-path-composition:3}

The \textbf{path algebra} \(kQ\) of a quiver \(Q\) over a field (or ring) \(k\) is the ring whose elements are \(k\)-linear combinations of paths in \(Q\), with multiplication given by path composition (composing two paths whose endpoints do not match gives \(0\)). Formalizing \(kQ\) fully (as a \texttt{Ring}, per Chapter 8) requires ``formal sums of paths with ring coefficients,'' which is a genuinely bigger construction: essentially a finitely-supported function from paths to \(k\). It is a natural next project once the material above is well understood. This chapter stops at ``paths and their composition'' because that data (the \emph{category} of paths, really) is the essential content of the construction; once it is in place, the ring structure on top is routine to add.

\begin{mathreading}

The \textbf{path algebra} \(kQ\) is the free \(k\)-module on the set of all paths, \(kQ = \bigoplus_{p\ \text{path}} k\cdot p\), with multiplication extending path composition \(k\)-bilinearly: \[
q \cdot p = \begin{cases} q\circ p & \text{if } t(p) = s(q),\\ 0 &
\text{otherwise,}\end{cases}
\] and unit \(\sum_{v\in V} e_v\) (the sum of the trivial paths). Equivalently \(kQ\) is the \(k\)-linearization of the free category \(\mathrm{Free}(Q)\): its \emph{category algebra}. So representations of \(Q\) are exactly \(kQ\)-modules, the bridge to Chapter 10 promised there. The construction requires finitely-supported functions \(\{\text{paths}\} \to k\), i.e.~the free module, which is why only the composition (the category layer) is formalized here.

\end{mathreading}

\begin{quote}
Read more: \hyperref[sec:10-modules:00-index:1]{Chapter 10}'s \href{https://loogle.lean-lang.org/?q=Module}{\texttt{Module}}, \href{https://loogle.lean-lang.org/?q=LinearMap}{\texttt{LinearMap}}, and direct-sum material is exactly the vocabulary a full \(kQ\)-module (i.e.~quiver representation) formalization would need. Externally, Assem--Simson--Skowroński's \emph{Elements of the Representation Theory of Associative Algebras} (Vol. 1) and Schiffler's \emph{Quiver Representations} both develop path algebras and their representations from scratch, at a similar level of explicitness to this chapter.
\end{quote}

\subsection{References}\label{sec:11-path-algebras:05-path-composition:4}

Full citations in the \hyperref[sec:root:bibliography:1]{Bibliography}.

\begin{itemize}
\tightlist
\item
  Assem, Simson, and Skowroński (\cite{AssemSimsonSkowronski2006}) --- the standard graduate reference for quivers, path algebras, and their representations, developed from scratch as in this chapter.
\item
  Schiffler (\cite{Schiffler2014}) --- an elementary, textbook-level treatment of the same material, pitched at advanced undergraduates/beginning graduate students.
\end{itemize}
